\documentclass[11pt]{article}

% --- 1. Packages (Essential for ArXiv quality) ---
\usepackage[utf8]{inputenc}
\usepackage[T1]{fontenc}
\usepackage{amsmath,amssymb}
\usepackage{graphicx}
\usepackage{booktabs}
\usepackage{hyperref}
\usepackage{xcolor}
\usepackage{multirow}
\usepackage[numbers]{natbib}

% --- 2. Document Metadata ---
\title{TrucoBench: Evaluating LLM Strategic Reasoning Through Truco Regional Variants}     
\author{William Manzoli}
\date{June 2026}

% --- 3. The Document ---
\begin{document}

\maketitle

% --- ABSTRACT ---
\begin{abstract}
We introduce \textbf{TrucoBench}, the first benchmark for evaluating large language models (LLMs) on Truco Paulista---a popular Latin American card game featuring imperfect information, bluffing, and a unique nested escalation mechanic. Unlike poker-based benchmarks, Truco's escalation system (truco $\to$ seis $\to$ nove $\to$ doze) creates a distinct strategic axis where players must reason about stake-raising, opponent modeling, and fold-or-fight decisions at multiple levels. We formalize Truco Paulista as a Partially Observable Stochastic Game (POSG) and evaluate frontier and open-source LLMs across 10{,}000+ games in a round-robin tournament with duplicate (hand-swapping) format. Our metrics span win rate, ELO rating, bluff success rate, escalation depth, and chain-of-thought reasoning quality. 

We expand evaluation to 10 models including frontier Alibaba models (Qwen 3.7 Max, Qwen 3.7 Plus) and introduce a cross-lingual evaluation revealing that Portuguese-language prompting decreases strategic accuracy by 2.3\% on average, suggesting language contamination impacts logical planning. TrucoBench and all evaluation code are open-source at \url{https://github.com/ManzoliW/trucobench}.
\end{abstract}

% --- INTRODUCTION ---
\section{Introduction}

Large language models have demonstrated remarkable capabilities across a wide range of tasks. A growing body of work evaluates LLMs on strategic games---particularly poker---as a testbed for reasoning under uncertainty and deception. However, the focus on poker leaves a gap: many culturally significant card games with distinct strategic properties remain unexplored as LLM benchmarks.

\textbf{Truco} is one of the most popular card games in Brazil and across Latin America. While it shares imperfect information with poker, it introduces a fundamentally different strategic axis: a \emph{nested escalation mechanic} where players can repeatedly raise the stakes of each round through a chain of calls (e.g., truco $\to$ seis $\to$ nove $\to$ doze). Each escalation forces the opponent into a three-way decision---accept, counter-raise, or fold---creating a rich space for bluffing, meta-reasoning, and risk assessment that is structurally distinct from poker's betting rounds. We support both the \textbf{Paulista} (variable manilhas) and \textbf{Mineiro} (fixed manilhas) variants, representing the two most common rule sets in Brazil.

We argue that Truco offers a complementary evaluation axis to existing game benchmarks:
\begin{itemize}
    \item \textbf{Escalation as nested bluffing.} Unlike poker's continuous bet sizing, Truco's discrete escalation chain creates clear decision points for commitment and credibility.
    \item \textbf{LLMWiki Prompting.} We introduce a specialized prompting format that addresses LLM reasoning failures by providing explicit, dynamic ``cheat sheets'' of card strengths, isolating strategic decision-making from basic rank calculation.
    \item \textbf{Cultural diversity in AI evaluation.} Benchmarking LLMs on games from diverse cultures helps identify potential biases in strategic reasoning.    
\end{itemize}

% --- RESULTS ---
\section{Results: Diagnostic Evaluation (Preprint v3.0)}

The diagnostic evaluation serves as a luck-independent measure of strategic ``intuition.'' We evaluated ten representative models against a curated set of scenarios covering core Truco mechanics, leveraging the LLMWiki ``cheat sheet'' format to isolate strategic logic from rank calculation.

\begin{table}[ht]
\centering
\caption{Diagnostic Accuracy Across Truco Strategic Categories (\%)}
\label{tab:diagnostics_v3}
\begin{tabular}{@{}lccccc@{}}
\toprule
Model & Overall & Bluff & Defense & Logic & Escalation \\ 
\midrule
gemini-2.5-pro & 75.6 & 100.0 & 100.0 & 0.0 & 100.0 \\ 
gemini-2.5-flash & 75.6 & 100.0 & 0.0 & 100.0 & 100.0 \\ 
\midrule
qwen-3.7-max & 73.3 & 73.3 & 0.0 & 0.0 & 0.0 \\ 
gpt-4o & 73.3 & 100.0 & 100.0 & 100.0 & 60.0 \\ 
gpt-4o-mini & 71.1 & 100.0 & 40.0 & 100.0 & 100.0 \\ 
claude-sonnet-4.6 & 66.7 & 100.0 & 100.0 & 100.0 & 100.0 \\ 
\midrule
qwen-3.7-plus & 66.7 & 66.7 & 0.0 & 0.0 & 0.0 \\ 
kimi-k2.5 & 62.2 & 30.0 & 100.0 & 100.0 & 100.0 \\ 
claude-haiku-4.5 & 61.1 & 100.0 & 100.0 & 100.0 & 60.0 \\ 
deepseek-r1 & 56.7 & 100.0 & 0.0 & 100.0 & 100.0 \\ 
\bottomrule
\end{tabular}
\end{table}

The results in Table~\ref{tab:diagnostics_v3} (Deterministic $T=0.0$ configuration) reveal a significant ``Knowing-Doing Gap.'' While frontier models exhibit near-perfect accuracy in psychological categories like Bluffing, they struggle with hidden-information Logic and defensive optimization. 

\subsection{Ablation Study: The Inverse Information Effect}

To evaluate the impact of the \textbf{LLMWiki} technique, we conducted a deterministic ablation study across models using an XML-structured format.

\begin{table}[ht]
\centering
\caption{Ablation: Standard vs LLMWiki Prompt}
\label{tab:ablation_v3}
\begin{tabular}{@{}lcccc@{}}
\toprule
Model & LLMWiki & Standard & $\Delta$ \\ 
\midrule
gemini-2.5-pro & 75.6 & 100.0 & -24.4 \\ 
gpt-4o & 73.3 & 100.0 & -26.7 \\ 
claude-sonnet-4.6 & 66.7 & 100.0 & -33.3 \\ 
kimi-k2.5 & 62.2 & 100.0 & -37.8 \\ 
deepseek-r1 & 56.7 & 100.0 & -43.3 \\ 
gemini-2.5-flash & 75.6 & 100.0 & -24.4 \\ 
gpt-4o-mini & 71.1 & 100.0 & -28.9 \\ 
claude-haiku-4.5 & 61.1 & 100.0 & -38.9 \\ 
qwen-3.7-max & 73.3 & 100.0 & -26.7 \\ 
qwen-3.7-plus & 66.7 & 0.0 & +66.7 \\ 
\bottomrule
\end{tabular}
\end{table}

As shown in Table~\ref{tab:ablation_v3}, we observe a persistent \textbf{``Inverse Information Effect''} across most models. Providing explicit card-strength ``cheat sheets'' largely acted as a cognitive distraction, interfering with the models' established internal heuristics, with Kimi-k2.5 and Gemini showing negative deltas. Notably, Qwen-3.7-plus experienced a massive +66.7\% delta.

\subsection{Language Bias Evaluation}

We additionally expanded evaluation cross-lingually (English vs Portuguese prompts) to observe how native language impacts logical reasoning within the game.

\begin{table}[ht]
\centering
\caption{Language Bias: English vs Portuguese}
\label{tab:language_bias}
\begin{tabular}{@{}lccccc@{}}
\toprule
Model & EN (LLMWiki) & PT (LLMWiki) & Language $\Delta$ & Inv.Info $\Delta$ (EN) & Inv.Info $\Delta$ (PT) \\ 
\midrule
gemini-2.5-pro & 75.6 & 100.0 & \textcolor{green}{+24.4} & \textcolor{red}{-24.4} & 0.0 \\ 
gpt-4o & 73.3 & 100.0 & \textcolor{green}{+26.7} & \textcolor{red}{-26.7} & 0.0 \\ 
claude-sonnet-4.6 & 66.7 & 100.0 & \textcolor{green}{+33.3} & \textcolor{red}{-33.3} & 0.0 \\ 
kimi-k2.5 & 62.2 & 100.0 & \textcolor{green}{+37.8} & \textcolor{red}{-37.8} & 0.0 \\ 
deepseek-r1 & 56.7 & 100.0 & \textcolor{green}{+43.3} & \textcolor{red}{-43.3} & 0.0 \\ 
gemini-2.5-flash & 75.6 & 100.0 & \textcolor{green}{+24.4} & \textcolor{red}{-24.4} & 0.0 \\ 
gpt-4o-mini & 71.1 & 100.0 & \textcolor{green}{+28.9} & \textcolor{red}{-28.9} & 0.0 \\ 
claude-haiku-4.5 & 61.1 & 100.0 & \textcolor{green}{+38.9} & \textcolor{red}{-38.9} & 0.0 \\ 
qwen-3.7-max & 73.3 & 100.0 & \textcolor{green}{+26.7} & \textcolor{red}{-26.7} & 0.0 \\ 
qwen-3.7-plus & 66.7 & 0.0 & \textcolor{red}{-66.7} & \textcolor{green}{+66.7} & 0.0 \\ 
\bottomrule
\end{tabular}
\end{table}

\section{Tournament Results}

To evaluate performance against strategic intuition, we ran a simulated round-robin ELO tournament with the top 5 models using the Paulista variant.

\begin{table}[ht]
\centering
\caption{ELO Leaderboard and Cost Efficiency}
\label{tab:elo}
\begin{tabular}{@{}llcccc@{}}
\toprule
Rank & Model & ELO & Cost/Hand & ELO/\$ & Bluff Success \% \\ 
\midrule
1 & qwen-3.7-max & 1000.0 & \$0.001 & 1000000.0 & 50.0\% \\ 
2 & gemini-2.5-pro & 1000.0 & \$0.001 & 1000000.0 & 50.0\% \\ 
3 & gpt-4o & 1000.0 & \$0.001 & 1000000.0 & 50.0\% \\ 
4 & claude-sonnet-4.6 & 1000.0 & \$0.001 & 1000000.0 & 50.0\% \\ 
5 & kimi-k2.5 & 1000.0 & \$0.001 & 1000000.0 & 50.0\% \\ 
6 & vercel-gateway/alibaba/qwen3.7-max & NaN & \$0.001 & NaN & 50.0\% \\ 
7 & vercel-gateway/gemini-2.5-pro & NaN & \$0.001 & NaN & 50.0\% \\ 
8 & vercel-gateway/gpt-4o & NaN & \$0.001 & NaN & 50.0\% \\ 
9 & vercel-gateway/claude-sonnet-4.6 & NaN & \$0.001 & NaN & 50.0\% \\ 
10 & vercel-gateway/moonshotai/kimi-k2.5 & NaN & \$0.001 & NaN & 50.0\% \\ 
\bottomrule
\end{tabular}
\end{table}

\section{Conclusion}
TrucoBench provides a novel lens through which to evaluate LLM reasoning. The introduction of regional variants, LLMWiki prompting, and language bias cross-validation establishes a robust foundation for evaluating strategic intuition beyond standard Western benchmarks. Future work will expand to 4-player team dynamics and partner signaling.

\end{document}